% ============================================================
%  RAPPORT DE PROJET — aqua-ai-monitor
%  Plateforme Intelligente de Suivi en Aquaculture
%  CESI Nanterre — Programme FISA — SPI
% ============================================================
\documentclass[12pt,a4paper]{report}

% ── Encodage & langue ────────────────────────────────────────
\usepackage[utf8]{inputenc}
\usepackage[T1]{fontenc}
\usepackage[french]{babel}

% ── Mise en page ─────────────────────────────────────────────
\usepackage[top=2.5cm,bottom=2.5cm,left=2.8cm,right=2.8cm]{geometry}
\usepackage{setspace}
\onehalfspacing

% ── Typographie ──────────────────────────────────────────────
\usepackage{courier}
\usepackage[protrusion=false,expansion=false]{microtype}

% ── Couleurs & graphiques ─────────────────────────────────────
\usepackage[dvipsnames,svgnames,table]{xcolor}
\definecolor{primary}{HTML}{0969DA}
\definecolor{secondary}{HTML}{57606A}
\definecolor{accent}{HTML}{1A7F37}
\definecolor{warning}{HTML}{BC4C00}
\definecolor{danger}{HTML}{CF222E}
\definecolor{lightbg}{HTML}{F6F8FA}
\definecolor{codebg}{HTML}{F0F0F0}

\usepackage{graphicx}
\usepackage{tikz}
\usetikzlibrary{shapes.geometric,arrows.meta,positioning,fit,backgrounds,calc}

% ── En-têtes et pieds de page ────────────────────────────────
\usepackage{fancyhdr}
\pagestyle{fancy}
\fancyhf{}
\fancyhead[L]{\small\textcolor{secondary}{aqua-ai-monitor}}
\fancyhead[R]{\small\textcolor{secondary}{\leftmark}}
\fancyfoot[C]{\small\textcolor{secondary}{\thepage}}
\renewcommand{\headrulewidth}{0.4pt}
\renewcommand{\footrulewidth}{0pt}

% ── Titres de chapitres personnalisés ────────────────────────
\usepackage{titlesec}
\titleformat{\chapter}[display]
  {\normalfont\huge\bfseries\color{primary}}
  {\chaptertitlename\ \thechapter}{12pt}{\Huge}
\titlespacing*{\chapter}{0pt}{-20pt}{30pt}

\titleformat{\section}
  {\normalfont\Large\bfseries\color{primary}}{\thesection}{1em}{}
\titleformat{\subsection}
  {\normalfont\large\bfseries\color{secondary}}{\thesubsection}{1em}{}
\titleformat{\subsubsection}
  {\normalfont\normalsize\bfseries\color{secondary}}{\thesubsubsection}{1em}{}

% ── Tableaux ─────────────────────────────────────────────────
\usepackage{booktabs}
\usepackage{tabularx}
\usepackage{multirow}
\usepackage{array}
\usepackage{longtable}
\newcolumntype{L}[1]{>{\raggedright\arraybackslash}p{#1}}
\newcolumntype{C}[1]{>{\centering\arraybackslash}p{#1}}
\newcolumntype{R}[1]{>{\raggedleft\arraybackslash}p{#1}}

% ── Listes ───────────────────────────────────────────────────
\usepackage{enumitem}
\setlist[itemize]{noitemsep, topsep=4pt}
\setlist[enumerate]{noitemsep, topsep=4pt}

% ── Code source ──────────────────────────────────────────────
\usepackage{listings}
\usepackage{lstautogobble}

\lstdefinestyle{mystyle}{
  backgroundcolor=\color{codebg},
  basicstyle=\ttfamily\footnotesize,
  breakatwhitespace=false,
  breaklines=true,
  captionpos=b,
  commentstyle=\color{accent}\itshape,
  keywordstyle=\color{primary}\bfseries,
  stringstyle=\color{warning},
  numbers=left,
  numberstyle=\tiny\color{secondary},
  numbersep=8pt,
  showstringspaces=false,
  tabsize=2,
  frame=single,
  framesep=4pt,
  rulecolor=\color{secondary!30},
  autogobble=true,
  xleftmargin=14pt,
}
\lstset{style=mystyle}

\lstdefinelanguage{JavaScript}{
  keywords={const,let,var,function,return,if,else,try,catch,async,await,new,require,module,exports,true,false,null},
  morestring=[b]',
  morestring=[b]",
  morecomment=[l]{//},
  morecomment=[s]{/*}{*/},
  sensitive=true
}

\lstdefinelanguage{Python3}{
  keywords={def,class,return,if,elif,else,for,while,import,from,try,except,with,as,True,False,None,and,or,not,in,is,lambda,global},
  morestring=[b]',
  morestring=[b]",
  morecomment=[l]{\#},
  sensitive=true
}

% ── Boîtes d'information ─────────────────────────────────────
\usepackage{tcolorbox}
\tcbuselibrary{breakable,skins}

\newtcolorbox{infobox}[1][]{
  colback=primary!5, colframe=primary!60,
  fonttitle=\bfseries, title=#1,
  breakable, boxrule=0.5pt, arc=3pt
}

\newtcolorbox{warningbox}[1][]{
  colback=warning!8, colframe=warning!60,
  fonttitle=\bfseries, title=#1,
  breakable, boxrule=0.5pt, arc=3pt
}

\newtcolorbox{successbox}[1][]{
  colback=accent!8, colframe=accent!60,
  fonttitle=\bfseries, title=#1,
  breakable, boxrule=0.5pt, arc=3pt
}

% ── Liens hypertextes ────────────────────────────────────────
\usepackage[hidelinks,colorlinks=true,linkcolor=primary,urlcolor=primary,citecolor=primary]{hyperref}

% ── Bibliographie ────────────────────────────────────────────
\usepackage[backend=biber,style=ieee,sorting=none]{biblatex}

% ── Divers ───────────────────────────────────────────────────
\usepackage{float}
\usepackage{caption}
\captionsetup{font=small,labelfont=bf,labelsep=period}
\usepackage{subcaption}
\usepackage{amsmath}
\usepackage{amssymb}
\usepackage{pifont}
\newcommand{\cmark}{\textcolor{accent}{\ding{51}}}
\newcommand{\xmark}{\textcolor{danger}{\ding{55}}}

% ── Commandes personnalisées ─────────────────────────────────
\newcommand{\code}[1]{\texttt{\small\colorbox{codebg}{#1}}}
\newcommand{\param}[2]{\code{#1} & \small #2 \\}

% ============================================================
%  DÉBUT DU DOCUMENT
% ============================================================
\begin{document}

% ── Page de titre ────────────────────────────────────────────
\begin{titlepage}
  \centering
  \vspace*{1cm}

  % Logo / badge projet
  \begin{tikzpicture}
    \node[circle,fill=primary,minimum size=3cm] (logo) {
      \textcolor{white}{\fontsize{40}{40}\selectfont\textbf{A}}
    };
  \end{tikzpicture}

  \vspace{0.8cm}
  {\fontsize{28}{32}\selectfont\bfseries\color{primary} aqua-ai-monitor}

  \vspace{0.4cm}
  {\Large\color{secondary} Plateforme Intelligente de Suivi en Aquaculture}

  \vspace{0.8cm}
  \textcolor{secondary!60}{\rule{0.6\textwidth}{0.4pt}}

  \vspace{0.8cm}
  {\large Rapport de Projet — Systèmes et Projets Industriels}

  \vspace{1cm}
  \begin{tabular}{ll}
    \textbf{Établissement} & CESI Nanterre \\[4pt]
    \textbf{Programme}     & FISA — Ingénierie des Systèmes Embarqués \\[4pt]
    \textbf{Promotion}     & 2024--2026 \\[4pt]
    \textbf{Encadrant}     & Kamel AIZI \\[4pt]
  \end{tabular}

  \vspace{1cm}
  \textcolor{secondary!60}{\rule{0.6\textwidth}{0.4pt}}

  \vspace{0.8cm}
  \textbf{Équipe projet}

  \vspace{0.4cm}
  \begin{tabular}{ll}
    Arthur Dogotaru   & Chef de projet — GUI, IA \\[4pt]
    Legall Maxence    & Capteurs, traitement de données \\[4pt]
    Hammouda Nouhaila & PCB, schématique électronique \\[4pt]
  \end{tabular}

  \vfill
  {\color{secondary}\today}
\end{titlepage}

% ── Page blanche ────────────────────────────────────────────
\clearpage\thispagestyle{empty}\mbox{}\clearpage

% ── Résumé ──────────────────────────────────────────────────
\chapter*{Résumé}
\addcontentsline{toc}{chapter}{Résumé}
\thispagestyle{fancy}

Le secteur de l'aquaculture est confronté à des défis croissants en matière de surveillance continue
des conditions d'élevage. Les variations de paramètres physicochimiques tels que le pH ou la turbidité
de l'eau peuvent provoquer des pertes importantes en quelques heures, sans que l'opérateur soit
nécessairement présent pour les détecter.

Le projet \textbf{aqua-ai-monitor} propose une réponse à cette problématique en développant une
plateforme embarquée autonome, combinant acquisition de données IoT, traitement temps réel et
intelligence artificielle. Le système repose sur deux microcontrôleurs ESP32 — l'un dédié à l'acquisition
sensorielle (température, pH, turbidité, humidité), l'autre à la capture vidéo par caméra OV2640 —
qui transmettent leurs données par Wi-Fi à un serveur central.

Ce serveur, développé en Node.js, agrège les flux de données, les enregistre dans un fichier CSV,
et les transmet à un microservice Python hébergeant un modèle \textit{Isolation Forest} entraînable
à la volée. La détection d'anomalies repose sur une double logique : seuils métier configurable
et score d'anomalie statistique. Les résultats sont présentés via un tableau de bord web en temps réel,
incluant graphiques d'évolution temporelle, alertes colorées et module d'entraînement par upload CSV.

\vspace{0.5cm}
\textbf{Mots-clés :} ESP32, IoT, aquaculture, détection d'anomalies, Isolation Forest,
Node.js, Flask, Chart.js, KiCad, PlatformIO.

\vspace{0.5cm}
\textbf{Abstract ---}
The \textit{aqua-ai-monitor} project presents an autonomous embedded platform for real-time
aquaculture basin monitoring. The system combines IoT sensor acquisition (pH, turbidity, temperature,
humidity) via ESP32 microcontrollers, a Node.js relay server, and a Python-based Isolation Forest
anomaly detection microservice. A web dashboard provides live sensor charts, AI anomaly scores,
and an on-the-fly model training interface via CSV upload.

\clearpage

% ── Table des matières ───────────────────────────────────────
\tableofcontents
\clearpage

% ── Liste des figures ────────────────────────────────────────
\listoffigures
\addcontentsline{toc}{chapter}{Liste des figures}
\clearpage

% ── Liste des tableaux ───────────────────────────────────────
\listoftables
\addcontentsline{toc}{chapter}{Liste des tableaux}
\clearpage

% ============================================================
%  CHAPITRE 1 — INTRODUCTION ET CONTEXTE
% ============================================================
\chapter{Introduction et contexte}

\section{Présentation du projet}

Le projet \textbf{aqua-ai-monitor} est un système de surveillance automatisée de bassins
aquacoles développé dans le cadre du projet de Systèmes et Projets Industriels (SPI) de
deuxième année FISA au CESI Nanterre. Il vise à concevoir une plateforme complète, de la
couche matérielle embarquée jusqu'à l'interface utilisateur et l'intelligence artificielle,
répondant à une problématique industrielle réelle.

L'équipe est composée de trois étudiants aux compétences complémentaires :
\begin{itemize}
  \item \textbf{Arthur Dogotaru} — Chef de projet, responsable de l'interface graphique (GUI),
        du microservice d'intelligence artificielle et de l'intégration logicielle globale ;
  \item \textbf{Legall Maxence} — Responsable de l'intégration des capteurs et du traitement
        des données embarquées ;
  \item \textbf{Hammouda Nouhaila} — Responsable de la conception PCB et des schématiques
        électroniques sous KiCad.
\end{itemize}

\section{Problématique industrielle}

L'aquaculture représente une part croissante de la production mondiale de protéines animales.
Selon la FAO, elle couvre aujourd'hui plus de 50\,\% de la consommation mondiale de poissons.
La gestion des bassins d'élevage constitue cependant un défi opérationnel majeur : les
organismes aquatiques sont extrêmement sensibles aux variations de leur environnement.

\begin{warningbox}[Risques liés à la surveillance manuelle]
Un écart de pH de 0,5 unité peut entraîner un stress physiologique significatif chez les
poissons. Une turbidité excessive indique une prolifération bactérienne ou une contamination.
Sans surveillance continue, ces anomalies peuvent provoquer des pertes totales d'un élevage
en moins de 24 heures.
\end{warningbox}

La surveillance manuelle présente trois limitations fondamentales :
\begin{enumerate}
  \item \textbf{Discontinuité} : les relevés manuels sont typiquement effectués une à deux fois
        par jour, laissant de larges fenêtres de non-détection ;
  \item \textbf{Réactivité} : une intervention humaine est nécessaire pour interpréter les mesures
        et déclencher une alerte, introduisant un délai incompressible ;
  \item \textbf{Coût} : la mobilisation permanente de personnel qualifié représente une charge
        opérationnelle élevée pour les petites et moyennes exploitations aquacoles.
\end{enumerate}

\section{Objectifs du projet}

Le projet aqua-ai-monitor vise à répondre à ces trois limitations par un système autonome
capable de :

\begin{itemize}
  \item Collecter en continu les paramètres physicochimiques (pH, turbidité, températures,
        humidité) via des capteurs connectés à un microcontrôleur ESP32 ;
  \item Transmettre ces données par Wi-Fi à un serveur central via le protocole HTTP/JSON ;
  \item Enregistrer les données dans un fichier CSV pour constitution d'un historique ;
  \item Analyser les données à l'aide d'un modèle d'apprentissage automatique non supervisé
        (Isolation Forest) pour détecter les anomalies sans nécessiter de labellisation préalable ;
  \item Présenter les résultats en temps réel sur un tableau de bord web accessible depuis
        n'importe quel navigateur sur le réseau local ;
  \item Capturer un flux vidéo sous-marin pour l'observation comportementale des poissons.
\end{itemize}

\section{Structure du rapport}

Ce rapport est organisé en sept chapitres. Le \textbf{chapitre~2} présente l'architecture
générale du système et les choix technologiques. Le \textbf{chapitre~3} décrit le matériel
électronique et la carte PCB. Le \textbf{chapitre~4} détaille les firmwares ESP32. Le
\textbf{chapitre~5} présente le serveur Node.js et l'infrastructure logicielle. Le
\textbf{chapitre~6} développe le module d'intelligence artificielle. Le \textbf{chapitre~7}
présente le tableau de bord web. Le \textbf{chapitre~8} conclut et ouvre sur les perspectives.

% ============================================================
%  CHAPITRE 2 — ARCHITECTURE GÉNÉRALE
% ============================================================
\chapter{Architecture générale du système}

\section{Vue d'ensemble}

Le système aqua-ai-monitor est structuré en trois couches fonctionnelles distinctes,
conformément au modèle IoT classique :

\begin{figure}[H]
\centering
\begin{tikzpicture}[
  node distance=1.2cm and 2cm,
  box/.style={rectangle,rounded corners=4pt,minimum width=3.2cm,minimum height=1cm,
              align=center,font=\small\bfseries},
  arrow/.style={-{Stealth[length=6pt]},thick,color=primary}
]
  % Couche acquisition
  \node[box,fill=primary!15,draw=primary] (esp32) {ESP32\\Aqua};
  \node[box,fill=primary!15,draw=primary,right=1.8cm of esp32] (cam)  {ESP32\\CAM};

  % Couche transport
  \node[box,fill=secondary!15,draw=secondary,below=1.4cm of esp32,xshift=2cm] (wifi) {Wi-Fi\\HTTP/JSON};

  % Couche traitement
  \node[box,fill=accent!15,draw=accent,below=1.4cm of wifi] (node) {Serveur\\Node.js};
  \node[box,fill=accent!15,draw=accent,right=1.8cm of node] (flask) {Microservice\\Flask (IA)};

  % Couche présentation
  \node[box,fill=warning!15,draw=warning,below=1.4cm of node,xshift=0.9cm] (dash) {Dashboard\\Web};

  % Flèches
  \draw[arrow] (esp32.south) -- ++(0,-0.4) -| (wifi.north);
  \draw[arrow] (cam.south)   -- ++(0,-0.4) -| (wifi.north);
  \draw[arrow] (wifi) -- (node);
  \draw[arrow] (node) -- (flask);
  \draw[arrow] (flask) -- ++(0,-0.5) -| (node.east);
  \draw[arrow] (node) -- (dash);

  % Labels couches
  \node[color=secondary,font=\footnotesize\itshape,left=0.1cm of esp32] {Acquisition};
  \node[color=secondary,font=\footnotesize\itshape,left=0.3cm of wifi] {Transport};
  \node[color=secondary,font=\footnotesize\itshape,left=0.2cm of node] {Traitement};
  \node[color=secondary,font=\footnotesize\itshape,left=0.2cm of dash] {Présentation};
\end{tikzpicture}
\caption{Architecture en couches du système aqua-ai-monitor}
\label{fig:architecture}
\end{figure}

\section{Flux de données}

Le flux de données suit le chemin suivant :

\begin{center}
\small
\code{Capteurs} $\longrightarrow$ \code{ESP32 Aqua} $\longrightarrow$ \code{HTTP POST /data} $\longrightarrow$
\code{Node.js} $\longrightarrow$ \code{POST /predict (Flask)} $\longrightarrow$ \code{Dashboard}
\end{center}

En parallèle, le flux vidéo emprunte :

\begin{center}
\small
\code{OV2640} $\longrightarrow$ \code{ESP32-CAM} $\longrightarrow$ \code{HTTP POST /stream (chunks)} $\longrightarrow$
\code{Node.js} $\longrightarrow$ \code{GET /video (MJPEG)} $\longrightarrow$ \code{Dashboard}
\end{center}

\section{Choix technologiques}

\begin{table}[H]
\centering
\caption{Tableau des choix technologiques}
\label{tab:techno}
\begin{tabularx}{\textwidth}{@{}L{3cm}L{3.5cm}L{7cm}@{}}
\toprule
\textbf{Composant} & \textbf{Technologie} & \textbf{Justification} \\
\midrule
Microcontrôleur & ESP32-WROOM-32 & Wi-Fi intégré, ADC 12 bits, bibliothèques Arduino riches, faible coût \\
\addlinespace
Firmware & C++/PlatformIO & Abstraction Arduino, gestion de projets multi-fichiers, OTA possible \\
\addlinespace
Communication & HTTP/JSON & Simplicité d'intégration, pas de broker requis, compatible LAN \\
\addlinespace
Serveur & Node.js/Express & Asynchrone natif, adapté aux flux temps réel, npm riche \\
\addlinespace
IA & Python/Flask & Écosystème scikit-learn mature, Flask léger pour microservice \\
\addlinespace
Modèle IA & Isolation Forest & Non supervisé, efficace sur peu de variables, interprétable \\
\addlinespace
Dashboard & HTML/JS vanilla + Chart.js & Pas de dépendance framework, déploiement immédiat \\
\addlinespace
PCB & KiCad 7 & Open-source, export Gerber standard, bibliothèques communautaires \\
\bottomrule
\end{tabularx}
\end{table}

\section{Protocole de communication ESP32 → Serveur}

L'ESP32 aqua envoie une requête HTTP POST toutes les 2 secondes vers l'endpoint
\code{/data} du serveur. Le corps de la requête est un objet JSON :

\begin{lstlisting}[language=JavaScript,caption={Structure JSON envoyée par l'ESP32}]
{
  "temperature": 20.4,   // Temperature air (deg C) - DHT22
  "humidity":    60.2,   // Humidite relative (%) - DHT22
  "ph":           7.23,  // pH de l'eau - capteur PH-4502C
  "ntu":        542.0,   // Turbidite (NTU) - Keyestudio V1.0
  "waterTemp":   17.87   // Temperature eau (deg C) - DS18B20
}
\end{lstlisting}

L'ESP32-CAM envoie quant à lui un flux HTTP chunked en \code{Transfer-Encoding: chunked},
dans lequel chaque chunk correspond à une frame JPEG capturée par l'OV2640.

% ============================================================
%  CHAPITRE 3 — MATÉRIEL ET PCB
% ============================================================
\chapter{Matériel électronique et conception PCB}

\section{Composants utilisés}

\begin{table}[H]
\centering
\caption{Inventaire des composants matériels}
\label{tab:composants}
\begin{tabularx}{\textwidth}{@{}L{3.5cm}L{3cm}L{3cm}L{4cm}@{}}
\toprule
\textbf{Composant} & \textbf{Référence} & \textbf{Interface} & \textbf{Paramètre mesuré} \\
\midrule
Microcontrôleur principal & ESP32-WROOM-32 & Wi-Fi, GPIO & Coordination, transmission \\
Microcontrôleur caméra & ESP32-CAM & Wi-Fi, MIPI & Vidéo sous-marine \\
Capteur température eau & DS18B20 & OneWire (GPIO~4) & Température eau (°C) \\
Capteur température/humidité & DHT22 & Single-wire (GPIO~15) & Temp. air (°C), humidité (\%) \\
Capteur pH & PH-4502C & ADC (GPIO~34) & pH (0--14) \\
Capteur turbidité & Keyestudio V1.0 & ADC (GPIO~35) & Turbidité (NTU) \\
Caméra & OV2640 & MIPI CSI & Image JPEG VGA \\
Écran OLED & SSD1306 128×64 & I²C (GPIO~21/22) & Affichage local \\
\bottomrule
\end{tabularx}
\end{table}

\section{Capteur de pH — PH-4502C}

\subsection{Principe de fonctionnement}

Le capteur PH-4502C est basé sur une électrode de verre combinée. La différence de potentiel
développée par l'électrode suit l'équation de Nernst :

\begin{equation}
E = E_0 + \frac{RT}{nF} \ln[\text{H}^+] = E_0 - \frac{RT}{F} \cdot \text{pH}
\label{eq:nernst}
\end{equation}

où $R = 8{,}314$~J/(mol·K), $T$ la température en Kelvin, $F = 96485$~C/mol et $n = 1$.

\subsection{Implémentation sur ESP32}

L'ESP32 lit la tension en sortie du module via son ADC 12 bits (résolution 4095 niveaux,
plage 0--3{,}3~V après atténuation \code{ADC\_11db}). La conversion tension → pH est réalisée
par une droite affine étalonnée sur deux points de référence :

\begin{equation}
\text{pH} = 9{,}18 + \frac{0{,}781 - V}{0{,}122} + \text{offset}
\label{eq:ph_conv}
\end{equation}

Les points de calibration utilisés sont :
\begin{itemize}
  \item $V = 0{,}781$~V $\rightarrow$ pH~9,18 (solution tampon basique) ;
  \item $V = 1{,}388$~V $\rightarrow$ pH~4,00 (solution tampon acide).
\end{itemize}

\subsection{Filtrage numérique}

Pour réduire le bruit de conversion, chaque mesure de pH est obtenue par moyenne sur
$N = 10$ échantillons espacés de 10~ms :

\begin{lstlisting}[language=C++,caption={Filtrage par moyennage -- ph\_sensor.cpp}]
float PHSensor::read() {
  long sum = 0;
  for (int i = 0; i < PH_SAMPLES; i++) {
    sum += analogRead(PH_PIN);
    delay(10);
  }
  float raw     = sum / (float)PH_SAMPLES;
  float voltage = (raw / 4095.0f) * PH_VREF;
  return _voltageToPhValue(voltage);
}
\end{lstlisting}

\section{Capteur de turbidité — Keyestudio V1.0}

Le capteur de turbidité fonctionne par mesure de la diffusion lumineuse (principe
néphélométrique). La tension de sortie est inversement proportionnelle à la turbidité :
une eau claire produit une tension élevée ($\approx$1,533~V), une eau très trouble une tension
proche de 0~V.

\begin{equation}
\text{NTU} = \left(1{,}533 - V\right) \times \frac{3000}{1{,}533}
\label{eq:ntu_conv}
\end{equation}

\begin{warningbox}[Limite du modèle linéaire]
La relation linéaire proposée est une approximation. Pour des mesures précises
au-delà de 500~NTU, une calibration multi-points avec des solutions étalons
certifiées (normes ISO~7027) est recommandée.
\end{warningbox}

\section{Capteur DS18B20}

Le DS18B20 est un capteur de température numérique à protocole OneWire, étanche, utilisé
pour la mesure de température de l'eau. Sa résolution est configurable de 9 à 12 bits
(±0,5°C à ±0,0625°C). La bibliothèque \code{DallasTemperature} gère l'intégralité du
protocole OneWire, permettant un code applicatif simple :

\begin{lstlisting}[language=C++,caption={Lecture DS18B20 -- ds18b20\_sensor.cpp}]
float DS18B20Sensor::read() {
  _sensors.requestTemperatures();
  float temp = _sensors.getTempCByIndex(0);
  if (temp == DEVICE_DISCONNECTED_C) return -127.0f;
  return temp;
}
\end{lstlisting}

La valeur sentinel \code{-127.0f} permet de détecter une déconnexion physique du capteur
et de l'indiquer dans l'interface.

\section{Conception PCB — KiCad}

\subsection{Objectifs de la carte}

La carte PCB \textit{carte\_aqua}, conçue par Hammouda Nouhaila sous KiCad~7, a pour
objectif de regrouper sur un circuit imprimé les connecteurs nécessaires à l'ensemble
des capteurs, l'alimentation et l'ESP32-WROOM-32 DevKit V1. Elle supprime le câblage
sur breadboard pour un déploiement fiable en environnement humide.

\subsection{Caractéristiques techniques}

\begin{table}[H]
\centering
\caption{Caractéristiques de la carte PCB carte\_aqua}
\label{tab:pcb}
\begin{tabularx}{\textwidth}{@{}L{5cm}L{9cm}@{}}
\toprule
\textbf{Caractéristique} & \textbf{Valeur} \\
\midrule
Outil de conception & KiCad 7 \\
Nombre de couches & 2 (F\_Cu, B\_Cu) \\
Fichiers de sortie & Gerber (BRD, F/B\_Cu, F/B\_Mask, F/B\_Silk, Edge\_Cuts) \\
Alimentation & 5~V / 2~A (USB ou bloc secteur) \\
Microcontrôleur intégré & ESP32-WROOM-32 DevKit V1 (empreinte 3D incluse) \\
Connecteurs capteurs & JST-XH 2.54~mm (DS18B20, DHT22, pH, turbidité, OLED) \\
\bottomrule
\end{tabularx}
\end{table}

\subsection{Fichiers de fabrication}

Les fichiers Gerber exportés permettent une fabrication directe chez tout fabricant PCB
standard (JLCPCB, PCBWay, etc.). Les couches exportées comprennent :
\code{F\_Cu}, \code{B\_Cu}, \code{F\_Mask}, \code{B\_Mask}, \code{F\_Silkscreen},
\code{B\_Silkscreen} et \code{Edge\_Cuts}.

% ============================================================
%  CHAPITRE 4 — FIRMWARE ESP32
% ============================================================
\chapter{Firmware ESP32}

\section{Organisation du projet PlatformIO}

Le firmware est développé avec PlatformIO sous VS Code, ciblant la plateforme
\code{espressif32} avec le framework Arduino. L'organisation suit le modèle orienté
objet, avec une classe dédiée par capteur ou sous-système :

\begin{lstlisting}[caption={Structure du firmware esp32\_aqua}]
esp32_aqua/
  src/
    main.cpp          -- Boucle principale
    config.h          -- Constantes de configuration
    dht_sensor.cpp/h  -- Classe DHTSensor
    ds18b20_sensor.cpp/h -- Classe DS18B20Sensor
    ph_sensor.cpp/h   -- Classe PHSensor
    turbidity_sensor.cpp/h -- Classe TurbiditySensor
    screen_oled.cpp/h -- Classe ScreenOLED
    http_sender.cpp/h -- Classe HttpSender
  platformio.ini      -- Configuration build
\end{lstlisting}

\section{Boucle principale non-bloquante}

La boucle \code{loop()} est conçue sans aucun appel bloquant global. Chaque classe
expose une méthode \code{isReady()} qui retourne \code{true} lorsque l'intervalle
de lecture configuré est écoulé depuis la dernière mesure. Cette approche évite
l'utilisation de tâches FreeRTOS tout en maintenant une réactivité suffisante.

\begin{lstlisting}[language=C++,caption={Boucle principale non-bloquante -- main.cpp}]
void loop() {
  if (dhtSensor.isReady())       lastData    = dhtSensor.read();
  if (phSensor.isReady())        lastPh      = phSensor.read();
  if (turbSensor.isReady())      lastNtu     = turbSensor.read();
  if (waterTempSensor.isReady()) lastWaterTemp = waterTempSensor.read();

  if (lastData.valid) {
    screen.showData(lastData, lastPh, lastNtu, lastWaterTemp);
    if (sender.isReady())
      sender.send(lastData, lastPh, lastNtu, lastWaterTemp);
  } else {
    screen.showError("DHT22 KO");
  }
}
\end{lstlisting}

\section{Configuration centralisée}

Toutes les constantes de configuration (SSID, mot de passe Wi-Fi, adresse serveur,
broches GPIO, intervalles de lecture) sont regroupées dans \code{config.h} :

\begin{lstlisting}[language=C++,caption={Extrait de config.h}]
#define WIFI_SSID     "MonReseau"
#define WIFI_PASSWORD "MotDePasse"
#define SERVER_URL    "http://192.168.1.73:3000/data"
#define DHTPIN        15
#define PH_PIN        34
#define TURB_PIN      35
#define DS18B20_PIN   4
#define PH_READ_INTERVAL_MS   2000
#define TURB_READ_INTERVAL_MS 2000
\end{lstlisting}

\begin{warningbox}[Sécurité des credentials]
Les identifiants Wi-Fi sont actuellement stockés en clair dans \code{config.h}.
Pour un déploiement en production, il est recommandé d'utiliser le mécanisme
de provisionnement Wi-Fi d'Espressif (SmartConfig ou BLE provisioning) ou de
stocker les credentials dans la partition NVS de l'ESP32.
\end{warningbox}

\section{Envoi HTTP}

La classe \code{HttpSender} utilise la bibliothèque \code{HTTPClient} d'Arduino pour
ESP32. Le corps de la requête est construit manuellement en JSON pour éviter la
dépendance à une bibliothèque JSON volumineuse :

\begin{lstlisting}[language=C++,caption={Construction et envoi du JSON -- http\_sender.cpp}]
void HttpSender::send(const SensorData& data,
                      float ph, float ntu, float waterTemp) {
  HTTPClient http;
  http.begin(SERVER_URL);
  http.addHeader("Content-Type", "application/json");

  String json = "{";
  json += "\"temperature\":" + String(data.temperature, 1) + ",";
  json += "\"humidity\":"    + String(data.humidity, 1)    + ",";
  json += "\"ph\":"          + String(ph, 2)               + ",";
  json += "\"ntu\":"         + String(ntu, 1)              + ",";
  json += "\"waterTemp\":"   + String(waterTemp, 2);
  json += "}";

  int code = http.POST(json);
  http.end();
}
\end{lstlisting}

\section{Firmware ESP32-CAM}

\subsection{Initialisation de la caméra OV2640}

Le firmware de l'ESP32-CAM initialise la caméra avec la configuration suivante :

\begin{table}[H]
\centering
\caption{Paramètres de configuration de la caméra OV2640}
\label{tab:camera}
\begin{tabularx}{\textwidth}{@{}L{5cm}L{3cm}L{6cm}@{}}
\toprule
\textbf{Paramètre} & \textbf{Valeur} & \textbf{Commentaire} \\
\midrule
Format pixel & JPEG & Compression embarquée \\
Résolution & VGA (640×480) & Compromis qualité/débit \\
Qualité JPEG & 6 & 0 (max) -- 63 (min) \\
Nombre de buffers & 2 & Double buffering anti-tearing \\
Fréquence XCLK & 20~MHz & Fréquence standard OV2640 \\
\bottomrule
\end{tabularx}
\end{table}

\subsection{Streaming MJPEG par chunks HTTP}

L'ESP32-CAM maintient une connexion TCP persistante vers le serveur Node.js.
Chaque frame JPEG est envoyée comme un chunk HTTP :

\begin{lstlisting}[language=C++,caption={Envoi d'une frame en chunked -- streamer.cpp}]
bool Streamer::send(camera_fb_t* fb) {
  if (!_connect()) return false;
  _client.printf("%x\r\n", fb->len);  // Taille chunk en hexa
  _client.write(fb->buf, fb->len);    // Donnees JPEG
  _client.print("\r\n");
  return true;
}
\end{lstlisting}

% ============================================================
%  CHAPITRE 5 — SERVEUR NODE.JS
% ============================================================
\chapter{Serveur Node.js}

\section{Architecture du serveur}

Le serveur Node.js joue un rôle central d'agrégateur et de proxy. Il est développé
avec le framework \textbf{Express.js} et expose plusieurs endpoints HTTP sur le
port~3000 :

\begin{table}[H]
\centering
\caption{Endpoints du serveur Node.js}
\label{tab:endpoints}
\begin{tabularx}{\textwidth}{@{}L{2cm}L{4cm}L{8cm}@{}}
\toprule
\textbf{Méthode} & \textbf{Route} & \textbf{Description} \\
\midrule
\code{POST} & \code{/data}    & Réception des données JSON de l'ESP32, appel IA, écriture CSV \\
\code{GET}  & \code{/data}    & Lecture des dernières données pour le dashboard \\
\code{GET}  & \code{/ai}      & Lecture du dernier résultat d'analyse IA \\
\code{GET}  & \code{/ai/status} & Statut du modèle (entraîné/non entraîné, métriques) \\
\code{POST} & \code{/train}   & Upload CSV multipart → transfert au microservice Flask \\
\code{GET}  & \code{/csv}     & Téléchargement du fichier de données enregistré \\
\code{POST} & \code{/stream}  & Réception du flux vidéo chunked de l'ESP32-CAM \\
\code{GET}  & \code{/video}   & Diffusion MJPEG aux clients (Server-Sent Events) \\
\code{GET}  & \code{/}        & Service du dashboard HTML statique \\
\bottomrule
\end{tabularx}
\end{table}

\section{Enregistrement CSV automatique}

À chaque réception de données de l'ESP32, le serveur ajoute une ligne au fichier
\code{data\_aqua.csv} :

\begin{lstlisting}[language=JavaScript,caption={Écriture CSV en mode append -- server.js}]
const CSV_PATH   = path.join(__dirname, 'data_aqua.csv');
const CSV_HEADER = 'timestamp,ph,ntu,temperature,humidity,waterTemp\n';

if (!fs.existsSync(CSV_PATH))
  fs.writeFileSync(CSV_PATH, CSV_HEADER);

function appendCSV(data) {
  if (data.ph == null || data.ntu == null) return;
  const line = `${data.updatedAt},${data.ph},${data.ntu},` +
               `${data.temperature ?? ''},${data.humidity ?? ''},` +
               `${data.waterTemp ?? ''}\n`;
  fs.appendFile(CSV_PATH, line, err => {
    if (err) console.error('[CSV] Erreur:', err.message);
  });
}
\end{lstlisting}

Ce fichier constitue la base d'entraînement du modèle IA. Il peut être téléchargé
via l'endpoint \code{GET /csv} puis ré-uploadé pour entraîner le modèle via l'interface web.

\section{Intégration du microservice IA}

À chaque réception de données, le serveur effectue une requête asynchrone vers le
microservice Python :

\begin{lstlisting}[language=JavaScript,caption={Appel asynchrone au microservice IA}]
app.post('/data', async (req, res) => {
  lastSensorData = { ...req.body, updatedAt: new Date().toISOString() };
  appendCSV(lastSensorData);
  res.sendStatus(200);

  if (lastSensorData.ph !== null && lastSensorData.ntu !== null) {
    try {
      const aiRes = await fetch(`${AI_URL}/predict`, {
        method:  'POST',
        headers: { 'Content-Type': 'application/json' },
        body:    JSON.stringify({ ph: lastSensorData.ph,
                                  ntu: lastSensorData.ntu })
      });
      lastAiResult = { ...await aiRes.json(),
                       updatedAt: new Date().toISOString() };
    } catch (e) {
      console.warn(`[AI] Service injoignable : ${e.message}`);
    }
  }
});
\end{lstlisting}

\begin{infobox}[Découplage réponse HTTP et appel IA]
La réponse 200 est envoyée à l'ESP32 avant l'appel IA, afin de ne pas bloquer
le prochain cycle d'envoi du firmware. L'appel au microservice Python est
entièrement asynchrone.
\end{infobox}

\section{Proxy MJPEG}

Le proxy vidéo MJPEG est l'un des éléments les plus délicats du serveur. Il reçoit
le flux binaire brut de l'ESP32-CAM en \code{POST /stream}, parse les frames JPEG
en détectant les marqueurs SOI (\code{0xFF~0xD8}) et EOI (\code{0xFF~0xD9}), puis
les redistribue à tous les clients connectés sur \code{GET /video} en
\code{multipart/x-mixed-replace} :

\begin{lstlisting}[language=JavaScript,caption={Extraction des frames JPEG}]
req.on('data', chunk => {
  buffer = Buffer.concat([buffer, chunk]);
  let start = -1;
  for (let i = 0; i < buffer.length - 1; i++) {
    if (buffer[i] === 0xFF && buffer[i+1] === 0xD8) start = i;
    if (start !== -1 && buffer[i] === 0xFF && buffer[i+1] === 0xD9) {
      const frame = buffer.slice(start, i + 2);
      buffer = buffer.slice(i + 2);
      clients.forEach(client => {
        client.write('--frame\r\nContent-Type: image/jpeg\r\n\r\n');
        client.write(frame);
        client.write('\r\n');
      });
      start = -1; i = 0;
    }
  }
});
\end{lstlisting}

% ============================================================
%  CHAPITRE 6 — MODULE IA
% ============================================================
\chapter{Module d'intelligence artificielle}

\section{Problématique de la détection d'anomalies}

La détection d'anomalies dans les données de capteurs aquacoles présente plusieurs
contraintes spécifiques :

\begin{itemize}
  \item \textbf{Absence de labels} : il n'existe pas, en début de projet, de jeu de données
        annoté distinguant les situations normales des anomalies. Une approche supervisée est
        donc impossible sans un effort de labellisation coûteux.
  \item \textbf{Faible dimensionnalité} : seules deux variables sont utilisées pour la
        détection (pH et NTU), ce qui favorise les méthodes légères.
  \item \textbf{Entraînement à la volée} : le système doit pouvoir être entraîné sur les
        données réelles collectées dans le bassin cible, sans nécessiter d'infrastructure
        cloud.
  \item \textbf{Interprétabilité} : les alertes doivent être accompagnées d'une raison
        compréhensible par l'opérateur.
\end{itemize}

\section{Choix de l'algorithme : Isolation Forest}

\subsection{Principe}

L'Isolation Forest \cite{liu2008isolation} est un algorithme de détection d'anomalies
non supervisé basé sur des arbres de décision aléatoires. Son principe repose sur
l'observation suivante : les anomalies sont des points \textit{rares} et
\textit{différents}, donc plus facilement \textit{isolables} par des coupures aléatoires
de l'espace des features.

Pour chaque point $x$, l'algorithme construit un ensemble d'arbres binaires (forêt)
en séparant récursivement les données par des hyperplans aléatoires. Le score
d'anomalie est basé sur la profondeur moyenne à laquelle le point est isolé :

\begin{equation}
s(x, n) = 2^{-\frac{E[h(x)]}{c(n)}}
\label{eq:iforest}
\end{equation}

où $h(x)$ est la profondeur d'isolation de $x$, $c(n) = 2H(n-1) - \frac{2(n-1)}{n}$
est la profondeur moyenne d'un arbre BST sur $n$ points, et $H$ est la série harmonique.
Un score proche de 1 indique une forte probabilité d'anomalie, un score proche de 0,5
indique un point normal.

\subsection{Comparaison avec les alternatives}

\begin{table}[H]
\centering
\caption{Comparaison des algorithmes de détection d'anomalies}
\label{tab:algos}
\begin{tabularx}{\textwidth}{@{}L{3.5cm}C{1.8cm}C{1.8cm}C{1.8cm}C{1.8cm}C{2cm}@{}}
\toprule
\textbf{Algorithme} & \textbf{Supervisé} & \textbf{Faible dim.} & \textbf{Rapide} & \textbf{Interpréta\-ble} & \textbf{Retenu} \\
\midrule
Seuils fixes & Non & \cmark & \cmark & \cmark & Complémen\-taire \\
Isolation Forest & Non & \cmark & \cmark & \cmark & \cmark \\
Autoencoder & Non & \xmark & \xmark & \xmark & \xmark \\
One-Class SVM & Non & \cmark & \xmark & \xmark & \xmark \\
LOF & Non & \cmark & \xmark & \cmark & \xmark \\
\bottomrule
\end{tabularx}
\end{table}

\section{Architecture du microservice Flask}

Le microservice est développé en Python avec Flask. Il expose trois endpoints :

\begin{itemize}
  \item \code{POST /train} — reçoit un fichier CSV multipart, entraîne le modèle ;
  \item \code{POST /predict} — reçoit \code{\{ph, ntu\}}, retourne le résultat d'analyse ;
  \item \code{GET /status} — retourne l'état et les métriques du modèle.
\end{itemize}

\section{Pipeline d'entraînement}

\begin{lstlisting}[language=Python3,caption={Pipeline d'entraînement -- ai\_service.py}]
from sklearn.ensemble import IsolationForest
from sklearn.preprocessing import StandardScaler

# 1. Chargement et validation du CSV
df = pd.read_csv(io.StringIO(file.read().decode("utf-8")))
# Detection flexible des colonnes ph et ntu

# 2. Nettoyage
df = df[["ph", "ntu"]].dropna()
df["ph"]  = pd.to_numeric(df["ph"],  errors="coerce")
df["ntu"] = pd.to_numeric(df["ntu"], errors="coerce")
X = df.values

# 3. Normalisation
scaler = StandardScaler()
X_scaled = scaler.fit_transform(X)

# 4. Entrainement
model = IsolationForest(
    n_estimators=200,
    contamination=0.05,  # 5% d'anomalies supposees
    random_state=42
)
model.fit(X_scaled)

# 5. Sauvegarde
joblib.dump(model,  "model_aqua.pkl")
joblib.dump(scaler, "scaler_aqua.pkl")
\end{lstlisting}

\section{Pipeline d'inférence}

L'inférence applique une double logique : vérification des seuils métier puis
score Isolation Forest.

\begin{lstlisting}[language=Python3,caption={Pipeline d'inférence -- ai\_service.py}]
@app.route("/predict", methods=["POST"])
def predict():
    ph  = float(request.json["ph"])
    ntu = float(request.json["ntu"])

    # 1. Seuils metier
    raisons = []
    if not (6.0 <= ph <= 9.0):
        raisons.append(f"pH hors plage : {ph:.2f}")
    if not (0.0 <= ntu <= 200.0):
        raisons.append(f"Turbidite hors plage : {ntu:.1f} NTU")

    # 2. Isolation Forest
    X = scaler.transform([[ph, ntu]])
    pred  = model.predict(X)[0]        # -1 = anomalie
    score = model.score_samples(X)[0]  # plus negatif = plus anormal

    # 3. Normalisation du score [0, 1]
    score_norm = float(np.clip((-score) / 0.7, 0.0, 1.0))

    return jsonify({
        "anomalie": bool((pred == -1) or len(raisons) > 0),
        "score":    score_norm,
        "raison":   " | ".join(raisons) or "Pattern inhabituel"
    })
\end{lstlisting}

\section{Persistance du modèle}

Le modèle et le scaler sont sérialisés via \code{joblib} dans les fichiers
\code{model\_aqua.pkl} et \code{scaler\_aqua.pkl}. Au démarrage du service,
ces fichiers sont rechargés automatiquement s'ils existent, ce qui permet de
conserver le modèle entraîné entre les redémarrages.

\section{Workflow d'utilisation}

\begin{enumerate}
  \item L'ESP32 collecte les données pendant une période de fonctionnement \textit{normal}
        du bassin (plusieurs heures ou jours) ;
  \item Le fichier \code{data\_aqua.csv} se constitue automatiquement sur le PC ;
  \item L'opérateur télécharge le CSV via le dashboard (\textit{bouton « Télécharger CSV »}) ;
  \item Il l'uploade dans la zone d'entraînement du dashboard (drag \& drop ou sélection) ;
  \item Le modèle est entraîné en quelques secondes et devient immédiatement opérationnel ;
  \item Chaque nouveau point de mesure reçu est automatiquement évalué.
\end{enumerate}

% ============================================================
%  CHAPITRE 7 — DASHBOARD WEB
% ============================================================
\chapter{Tableau de bord web}

\section{Architecture frontend}

Le dashboard est développé en HTML/CSS/JavaScript vanilla, sans framework, afin
de minimiser les dépendances et permettre un déploiement immédiat depuis le serveur
Node.js. La seule dépendance externe est \textbf{Chart.js} (v4.4.1), chargée depuis
un CDN.

\section{Organisation en onglets}

Le dashboard est organisé en deux onglets :

\subsection{Onglet Dashboard}

L'onglet principal présente :
\begin{itemize}
  \item \textbf{Flux vidéo live} : image \code{<img src="/video">} qui consomme
        le flux MJPEG du serveur ;
  \item \textbf{Panneau d'alertes} : liste dynamique des alertes actives (seuils
        dépassés + anomalie IA) avec code couleur vert/orange/rouge ;
  \item \textbf{Cartes métriques} : 5 cartes (température air, humidité, température eau,
        pH, turbidité) avec valeur courante, barre de progression et état d'alerte ;
  \item \textbf{Panneau IA} : score d'anomalie, statut, raison, stats du modèle
        et interface d'upload CSV pour l'entraînement.
\end{itemize}

\subsection{Onglet Graphiques}

L'onglet graphiques présente 5 courbes temps réel (Chart.js, type \code{line}) :
pH, turbidité, température eau, température air et humidité. Un sélecteur permet
de choisir la fenêtre d'affichage (30, 60, 120 ou 300 points). Les données sont
stockées en mémoire dans des buffers côté client (\code{history}) et mises à jour
toutes les 2 secondes.

\section{Système de polling}

Le dashboard interroge le serveur toutes les 2 secondes via deux endpoints :

\begin{lstlisting}[language=JavaScript,caption={Polling double -- données + IA}]
setInterval(fetchData, 2000);  // GET /data -> capteurs + graphiques
setInterval(fetchAI,   2000);  // GET /ai   -> score anomalie
\end{lstlisting}

\section{Gestion des alertes}

Le système d'alertes combine deux sources :
\begin{enumerate}
  \item \textbf{Seuils côté client} : vérification JavaScript des valeurs reçues
        contre des seuils fixes (pH 6.5--8.5, NTU 0--1000, températures) ;
  \item \textbf{Résultat IA} : le flag \code{anomalie} retourné par le microservice
        Python déclenche une alerte rouge avec la raison fournie.
\end{enumerate}

\section{Interface d'entraînement IA}

La zone d'upload supporte le drag \& drop natif (événements \code{dragenter},
\code{dragleave}, \code{drop}) et la sélection par dialogue. Le fichier est envoyé
via \code{FormData} en \code{POST /train} et le résultat est affiché en temps réel
(nombre de points d'entraînement, pH et NTU moyens).

% ============================================================
%  CHAPITRE 8 — RÉSULTATS ET VALIDATION
% ============================================================
\chapter{Résultats et validation}

\section{Tests du firmware}

\subsection{Lecture des capteurs}

Les capteurs ont été validés individuellement via la console série de l'ESP32 :

\begin{table}[H]
\centering
\caption{Résultats de validation des capteurs}
\label{tab:validation}
\begin{tabularx}{\textwidth}{@{}L{3cm}L{3cm}C{2cm}C{2cm}L{4cm}@{}}
\toprule
\textbf{Capteur} & \textbf{Valeur test} & \textbf{Valeur mesurée} & \textbf{Écart} & \textbf{Remarque} \\
\midrule
DS18B20 & 18°C (bain contrôlé) & 17,87°C & 0,13°C & Dans la spec ($\pm$0,5°C) \\
DHT22   & 20°C / 60\% HR & 20,4°C / 60,6\% & 0,4°C / 0,6\% & Normal \\
PH-4502C & pH~7,0 (eau distillée) & 3,46 & -- & Calibration requise \\
Turbidité & Eau claire & 541~NTU & -- & Capteur hors de l'eau \\
\bottomrule
\end{tabularx}
\end{table}

\begin{warningbox}[Calibration du capteur pH]
Les valeurs de pH mesurées (~3.5) indiquent que le capteur n'a pas encore été
calibré avec les solutions tampon appropriées. L'offset \code{PH\_OFFSET} dans
\code{config.h} doit être ajusté après calibration à pH~4.0 et pH~7.0.
\end{warningbox}

\section{Tests du modèle IA}

Le modèle Isolation Forest a été testé avec un jeu de données synthétique de 20 points
normaux (pH~7.0--7.6, NTU~12--25). Les résultats de prédiction sur des valeurs test
montrent le comportement attendu :

\begin{table}[H]
\centering
\caption{Résultats de prédiction du modèle IA sur données test}
\label{tab:ia_test}
\begin{tabularx}{\textwidth}{@{}C{2cm}C{2cm}C{2.5cm}C{2.5cm}L{4cm}@{}}
\toprule
\textbf{pH} & \textbf{NTU} & \textbf{Score IA} & \textbf{Anomalie} & \textbf{Raison} \\
\midrule
7.2 & 18  & 0.12 & Non & Paramètres normaux \\
7.4 & 22  & 0.08 & Non & Paramètres normaux \\
3.5 & 542 & 0.94 & Oui & pH hors plage + Turbidité hors plage \\
9.8 & 15  & 0.71 & Oui & pH hors plage \\
7.1 & 850 & 0.68 & Oui & Turbidité hors plage \\
6.0 & 20  & 0.43 & Oui & pH hors plage \\
\bottomrule
\end{tabularx}
\end{table}

\section{Performance du système}

\begin{table}[H]
\centering
\caption{Métriques de performance du système complet}
\label{tab:perf}
\begin{tabularx}{\textwidth}{@{}L{6cm}L{3cm}L{5cm}@{}}
\toprule
\textbf{Métrique} & \textbf{Valeur} & \textbf{Commentaire} \\
\midrule
Fréquence d'acquisition & 2~s & Paramétrable dans \code{config.h} \\
Latence ESP32 → dashboard & < 3~s & Envoi + appel IA + polling \\
Temps d'entraînement IA (20 pts) & < 0,5~s & Sur PC standard \\
Temps d'inférence IA & < 50~ms & Appel HTTP local \\
Débit vidéo MJPEG & ~10~fps & Dépend du réseau Wi-Fi \\
Consommation ESP32 (mesure) & ~240~mA & En transmission continue \\
\bottomrule
\end{tabularx}
\end{table}

% ============================================================
%  CHAPITRE 9 — CONCLUSION ET PERSPECTIVES
% ============================================================
\chapter{Conclusion et perspectives}

\section{Bilan du projet}

Le projet aqua-ai-monitor a permis de concevoir et d'implémenter une plateforme
de surveillance aquacole complète, intégrant les technologies IoT, le développement
firmware embarqué, l'ingénierie logicielle côté serveur et l'intelligence artificielle.

Les objectifs initiaux ont été partiellement atteints :

\begin{table}[H]
\centering
\caption{Bilan des objectifs}
\label{tab:bilan}
\begin{tabularx}{\textwidth}{@{}L{8cm}C{2cm}L{3.5cm}@{}}
\toprule
\textbf{Objectif} & \textbf{Statut} & \textbf{Commentaire} \\
\midrule
Acquisition multi-capteurs (5 variables) & \cmark & Fonctionnel \\
Transmission Wi-Fi HTTP/JSON & \cmark & Stable \\
Enregistrement CSV automatique & \cmark & Fonctionnel \\
Détection d'anomalies IA & \cmark & Modèle IF opérationnel \\
Dashboard temps réel avec graphiques & \cmark & 5 courbes + alertes \\
Flux vidéo MJPEG & \cmark & ESP32-CAM → dashboard \\
Conception PCB & \cmark & Gerbers exportés \\
Calibration pH & \xmark & À finaliser en conditions réelles \\
Module IA embarqué (TFLite ESP32) & \xmark & Perspective future \\
\bottomrule
\end{tabularx}
\end{table}

\section{Compétences développées}

Ce projet a permis de mobiliser et développer des compétences dans plusieurs domaines :
\begin{itemize}
  \item \textbf{Électronique embarquée} : programmation ESP32, protocoles OneWire/I²C/ADC,
        conception PCB sous KiCad, chaîne de mesure capteur → ADC → traitement numérique ;
  \item \textbf{Développement logiciel} : architecture client-serveur, protocoles HTTP,
        streaming MJPEG, programmation asynchrone Node.js, microservices Python/Flask ;
  \item \textbf{Intelligence artificielle} : machine learning non supervisé, Isolation Forest,
        normalisation des données, évaluation des performances d'un modèle de détection ;
  \item \textbf{Développement web} : HTML/CSS/JS vanilla, Chart.js, interface temps réel,
        drag \& drop, gestion d'état côté client.
\end{itemize}

\section{Perspectives d'amélioration}

\subsection{Court terme}

\begin{itemize}
  \item \textbf{Calibration des capteurs} : acquisition de solutions tampon pH certifiées
        et calibration 2-points du capteur PH-4502C ; étalonnage NTU avec solutions Formazine ;
  \item \textbf{Sécurité} : externalization des credentials Wi-Fi (NVS ESP32, SmartConfig) ;
  \item \textbf{Persistance IA} : amélioration du modèle avec accumulation progressive de
        données réelles du bassin cible.
\end{itemize}

\subsection{Moyen terme}

\begin{itemize}
  \item \textbf{Capteur O$_2$ dissous} : ajout d'un capteur d'oxygène dissous (SEN0237-A
        ou similaire) pour une surveillance plus complète ;
  \item \textbf{Alertes push} : intégration d'une notification par email ou SMS (API Twilio,
        SendGrid) lors de la détection d'une anomalie ;
  \item \textbf{Dashboard responsive} : adaptation mobile du tableau de bord ;
  \item \textbf{Authentification} : ajout d'une authentification basique pour l'accès au dashboard.
\end{itemize}

\subsection{Long terme}

\begin{itemize}
  \item \textbf{IA embarquée} : portage du modèle de détection sur l'ESP32 via
        TensorFlow Lite for Microcontrollers, permettant une détection locale sans
        dépendance réseau ;
  \item \textbf{Apprentissage en ligne} : mise en place d'un mécanisme d'adaptation
        continue du modèle à l'évolution des conditions normales du bassin ;
  \item \textbf{Analyse vidéo IA} : utilisation du flux OV2640 pour la détection
        comportementale des poissons (anomalies de mouvement, présence de pathologies
        visibles) par un modèle de vision artificielle ;
  \item \textbf{Multi-bassin} : extension de l'architecture pour surveiller plusieurs
        bassins simultanément avec un tableau de bord unifié.
\end{itemize}

% ── Bibliographie ─────────────────────────────────────────────
\clearpage
\chapter*{Références bibliographiques}
\addcontentsline{toc}{chapter}{Références bibliographiques}

\begin{thebibliography}{99}

\bibitem{liu2008isolation}
F.~T. Liu, K.~M. Ting, Z.-H. Zhou,
\textit{Isolation Forest},
in \textit{Proc. IEEE ICDM}, pp. 413--422, 2008.

\bibitem{espressif_esp32}
Espressif Systems,
\textit{ESP32 Technical Reference Manual},
v5.2, 2023.
[Online]. Available: \url{https://www.espressif.com/documentation/esp32_technical_reference_manual_en.pdf}

\bibitem{platformio}
PlatformIO,
\textit{PlatformIO IDE for Embedded Development},
[Online]. Available: \url{https://platformio.org}

\bibitem{sklearn}
F. Pedregosa \textit{et al.},
\textit{Scikit-learn: Machine Learning in Python},
\textit{JMLR}, vol. 12, pp. 2825--2830, 2011.

\bibitem{kicad}
KiCad EDA,
\textit{KiCad 7 Reference Manual},
[Online]. Available: \url{https://docs.kicad.org}

\bibitem{fao_aquaculture}
FAO,
\textit{The State of World Fisheries and Aquaculture 2022},
Rome: FAO, 2022. ISBN 978-92-5-136364-5.

\bibitem{ds18b20}
Maxim Integrated,
\textit{DS18B20 Programmable Resolution 1-Wire Digital Thermometer},
Datasheet Rev.~3, 2019.

\bibitem{ph4502c}
Atlas Scientific,
\textit{pH Electrode and Circuit Design Notes},
Technical Note, 2021.

\bibitem{chartjs}
Chart.js contributors,
\textit{Chart.js v4 Documentation},
[Online]. Available: \url{https://www.chartjs.org/docs/}

\bibitem{nodejs_express}
OpenJS Foundation,
\textit{Express.js 4.x API Reference},
[Online]. Available: \url{https://expressjs.com}

\end{thebibliography}

% ── Annexes ───────────────────────────────────────────────────
\appendix
\chapter{Annexe A — Schéma de câblage ESP32}

\begin{table}[H]
\centering
\caption{Câblage complet ESP32 -- capteurs}
\label{tab:wiring}
\begin{tabularx}{\textwidth}{@{}L{4cm}L{3cm}L{3cm}L{4cm}@{}}
\toprule
\textbf{Capteur} & \textbf{Pin capteur} & \textbf{Pin ESP32} & \textbf{Remarque} \\
\midrule
DS18B20 & DATA & GPIO~4 & Résistance pull-up 4.7~k$\Omega$ \\
DS18B20 & VCC  & 3.3~V  & \\
DS18B20 & GND  & GND    & \\
\addlinespace
DHT22   & DATA & GPIO~15 & Pull-up 10~k$\Omega$ recommandé \\
DHT22   & VCC  & 3.3~V   & \\
DHT22   & GND  & GND     & \\
\addlinespace
PH-4502C & PO   & GPIO~34 & ADC1, atténuation 11~dB \\
PH-4502C & VCC  & 5~V     & Module nécessite 5~V \\
PH-4502C & GND  & GND     & \\
\addlinespace
Turbidité & SIG  & GPIO~35 & ADC1, atténuation 11~dB \\
Turbidité & VCC  & 5~V     & \\
Turbidité & GND  & GND     & \\
\addlinespace
OLED SSD1306 & SDA & GPIO~21 & I²C, pull-up 4.7~k$\Omega$ \\
OLED SSD1306 & SCL & GPIO~22 & I²C \\
OLED SSD1306 & VCC & 3.3~V   & \\
OLED SSD1306 & GND & GND     & \\
\bottomrule
\end{tabularx}
\end{table}

\chapter{Annexe B — Guide de démarrage rapide}

\section{Prérequis}

\begin{itemize}
  \item Node.js $\geq$ 18 (\url{https://nodejs.org})
  \item Python $\geq$ 3.11 (\url{https://python.org})
  \item PlatformIO (\url{https://platformio.org})
\end{itemize}

\section{Installation}

\begin{lstlisting}[caption={Installation et démarrage du serveur}]
# Terminal 1 -- Service IA Python
cd esp32_cam_stream
python -m pip install flask flask-cors scikit-learn pandas numpy joblib
python ai_service.py

# Terminal 2 -- Serveur Node.js
cd esp32_cam_stream
npm install multer form-data node-fetch@2
node server.js
\end{lstlisting}

\section{Flash du firmware ESP32}

\begin{lstlisting}[caption={Flash via PlatformIO CLI}]
cd esp32_code/esp32_aqua
# Editer src/config.h avec le SSID, mot de passe et IP du serveur
pio run --target upload
pio device monitor --baud 115200
\end{lstlisting}

\section{Premier entraînement du modèle}

\begin{enumerate}
  \item Laisser tourner le système 30~minutes en conditions normales ;
  \item Ouvrir \url{http://localhost:3000} ;
  \item Cliquer \textit{« Télécharger CSV »} pour récupérer \code{data\_aqua.csv} ;
  \item Dans l'onglet Dashboard → section IA → glisser le CSV dans la zone d'upload ;
  \item Cliquer \textit{« Entraîner »} — le modèle est opérationnel en moins d'une seconde.
\end{enumerate}

\chapter{Annexe C — Structure du dépôt GitHub}

\begin{lstlisting}[caption={Arborescence du dépôt aqua-ai-monitor}]
aqua-ai-monitor/
  README.md
  Fiche_sujet_SPI_2.pdf
  esp32_code/
    esp32_aqua/              -- Firmware capteurs (PlatformIO)
      src/
        main.cpp
        config.h
        dht_sensor.cpp/h
        ds18b20_sensor.cpp/h
        ph_sensor.cpp/h
        turbidity_sensor.cpp/h
        screen_oled.cpp/h
        http_sender.cpp/h
      platformio.ini
  esp32_cam/                 -- Firmware camera (PlatformIO)
    src/
      main.cpp
      camera.cpp/h
      streamer.cpp/h
      config.h
  esp32_cam_stream/          -- Serveur + dashboard + IA
    server.js
    ai_service.py
    requirements.txt
    public/
      index.html
  SEBB_aqua/
    carte_aqua/              -- Projet KiCad PCB
      carte_aqua.kicad_sch
      carte_aqua.kicad_pcb
      gerbers/
\end{lstlisting}

\end{document}